\section{A Reusable Protocol for Responses and Paired Task Outcomes}
\label{sec:supp-protocol-card}
How could another study run the paired evaluation of main Secs.~\ref*{M-sec:comparisons} and~\ref*{M-sec:motion}, and what would it report? Each perturbed run is compared with the clean run that shares its policy noise, and the two clean runs give the clean rerun, the variation expected without any perturbation. Motion, lost and gained successes and the success change are reported against it, with bounds for missing outcomes. Section~\ref{sec:supp-protocol-validation} recomputes our estimates from the recorded case-level data, which needs no policy (producing those data does), and lists our implementation checks.

\noindent\textbf{1. Fix the plan before observing outcomes,} so that nothing is chosen by its result. Declare the model and checkpoint, task and visual variant, matched clean rendering, instruction, initial state, two policy-noise draws (0 and 1), episode budget and source code. Run all runs of a setting on one machine, using the five runs below in this order. Record missing or invalid runs against the plan. Policy, execution and setting together identify a comparison.

\begin{table}[!htb]\centering\small
\begin{tabular}{lll}
\toprule Run & Visual input / noise & Purpose\\\midrule
$R_0$ & matched clean / 0 & clean run\\
$T_0$ & perturbed / 0 & perturbed, noise of $R_0$\\
$R_1$ & matched clean / 1 & clean rerun\\
$T_1$ & perturbed / 1 & perturbed, noise of $R_1$\\
$R'_0$ & matched clean / 0 & repeat of $R_0$\\
\bottomrule
\end{tabular}
\end{table}

\noindent\textbf{2. Record plans, commands and motion separately.} Record the policy's raw actions, the planned commands after the output adapter, the commands sent to the environment and each call's return. Differences in the first five commands (translation, rotation, gripper) are in adapter command units; their root mean square (RMS) gives the size of a change. Record the end-effector position (world meters) and orientation (radians, with the quaternion convention) after returned controls 0, 5, 15, 30 and 60. Record each run's final pose, its step index and whether the two runs of a pair have the same number of steps, and early stops, unreturned controls, missing poses and unknown outcomes $Y$ explicitly. Fixed-step and final-pose comparisons use the recorded pose at that step. A returned command records what was sent to the controller; the signal it physically applied requires separate logging.

\noindent\textbf{3. Compare the perturbation with the clean rerun.} For each noise draw, compare $R_n$ with $T_n$. $R_0$--$R'_0$ checks reproducibility. $R_0$--$R_1$ is the clean rerun, conditional on the two recorded draws, and $T_0$--$T_1$ is its perturbed counterpart. For valid binary outcomes, SF is a lost success ($Y_R=1,Y_T=0$) and FS a gained success. The success change is $\Pr(\mathrm{FS})-\Pr(\mathrm{SF})$, and their sum is the rate of changed outcomes. Weight the two noise draws one half each, settings equally within a task, and tasks equally. Report bounds over missing outcomes and whole-task bootstrap intervals, adjusted for four families where stated, for sampling uncertainty.

\noindent\textbf{4. State where each internal quantity is recorded,} and keep recorded and recomputed states separate. For optional internal recordings of $\pi_{0.5}$, state for every scalar its layer, token region, generation step and recording path. Recomputing the prefix records the pre-block state $H$, the normalized state $Z$ and the projected cached values $V$ (fixed within a query). Recomputing the attention records, at the first query and generation step, the full-key probability $P$, the output after projection and before the gate, $W$, and the block-end action-expert state $H$. Keep L1, L9 and L18 and the main-image, wrist-image and text regions where available, and state whether a measure is raw or relative. Reconstruction coefficients describe an observed response; prediction and interventions need their own recorded inputs and operations. The $\pi_0$ paired executions stored no internal states; Sec.~\ref{sec:supp-pi0-internals} regenerates first-query states for 369 pairs, and the protocol evidence for $\pi_0$ is commands, motion at fixed steps and outcomes $Y$.

\noindent\textbf{5. Report task cost on its matched support.} For every comparison with the clean run, perturbed or clean rerun, state the settings, metric and time step it covers, including missing and undefined measurements, and pair the two means on their common cases. Post hoc SF/SS groups describe lost and kept successes among observed outcomes. Later data collections keep their own identities beside the original results, including negative, null and incomplete comparisons.

\subsection{Recomputing estimates from recorded case data}
\label{sec:supp-protocol-validation}
From the 536 policy--setting records, the analysis recomputes task-equal outcome estimates, bounds over missing outcomes and end-effector position and orientation comparisons after 5 and 60 controls. It checks that records are unique and complete, with five outcome cells per case and valid units and measurement status. It keeps the original task order and resampling scheme. The primary population has 202 planned settings per policy; a sensitivity analysis uses all 268.

Outcome estimates keep the planned denominator. For each case, all possible assignments of the unknown outcomes are enumerated before each comparison is formed. Each reported bound holds for its comparison alone; different comparisons can reach their endpoints under different assignments. Motion estimates require both noise-matched comparisons and the clean rerun at the same metric and time step. We average over noise draws, settings and tasks in that order, then divide the two task-equal means; a zero clean-rerun mean leaves the ratio undefined. Statistics: the success-change intervals use four-family-adjusted 98.75\% coverage, and other outcome and motion intervals are pointwise 95\%. Bounds over missing outcomes describe what the missing data allow, not sampling uncertainty.

\paragraph{Implementation and verification.}
Each entry names what was checked and what agreed, on the stated inputs rather than the original executions; removal-test and $\pi_0$ entries also record data-collection exceptions.
\begin{description}\small
\item[Recomputed hidden states and velocities.] One real input is run through two identical branches at generation steps 0, 5 and 9. The recomputed hidden states agree with the recorded ones; the recomputed velocity differs from the policy's by at most $0.03125$ in absolute value, with relative errors $0.00367$--$0.00457$. With identical branches, an error relative to the effect is undefined. $P$ and $W$ were not recorded during the policy's own generation, so attention has no tolerance or pass rate (Table~\ref{tab:recording-roles}).
\item[Case-level analysis.] Run from its instructions in an empty directory, the outcome and motion analysis reproduces 16 outcome panels and 128 metric/time/reference panels and passes 1,448 checks against the original values. The package includes the data dictionary, configuration, support tables, expected outputs and edge-case tests.
\item[Policies.] The configurations are \texttt{pi05\_libero} and \texttt{pi0\_libero}. Checksums of the checkpoint, normalization and tokenizer files of both policies are recorded in the source package.
\item[Reconstruction.] The uncentered 16-direction reconstruction results are reproduced to within $2.88\times10^{-14}$ in absolute task-score difference.
\item[Affine readout.] The implementation reproduces 324 recorded full-panel MSE and target-energy entries with the uncentered directions, with maximum absolute discrepancy zero.
\item[Nonlinear readouts.] The affine readout inside the nonlinear-readout code (Sec.~\ref{sec:supp-readout-capacity}) reproduces the expected affine results with absolute differences below $10^{-15}$ in task-equal MSE and $2\times10^{-13}$ in unscaled relative error reduction, in all four families.
\item[Corrected readout support.] The reported readouts use the protocol's 404 pairs with matched commands. A first run had evaluated Light on all 110 common-task pairs instead of the specified 96 with matched commands, which moved its affine values by up to $0.039$ in relative error reduction. That first run is kept as a record; its other three families are unchanged.
\item[Later-observation summaries.] The shared, clean-feature and random summaries in Table~\ref{tab:p08-dynamic-ranges} use the original averaging over noise, initial state and task and its intervals. Applying this averaging to all 940 available L18 episodes reproduces all 480 task means and 30 condition means.
\item[Aggregation.] The aggregation estimates and correlation intervals of the ten-condition study (Sec.~\ref{supp:legacy-d3}) are reproduced from the recorded case-level scalars.
\item[Removal test: edit size.] In the removal test (Sec.~\ref{sec:supp-native-ablation}), edit sizes are matched on the training average over the three conditions. Per condition, the differences exceed the 2\% tolerance.
\item[Removal test: data collection.] Data collection was paused at 229 of 288 cases and resumed in the fixed dispatch order. No case was selected by outcome. Two interim summaries (at 120 and 220 cases) had been seen before the decision to resume, as the decision record in the package states.
\item[Removal test: execution exceptions.] One case was interrupted after initialization; another was replaced under the prespecified technical-failure rule after a startup synchronization failure. Both are identified in the execution records.
\item[Removal test: numerical checks.] An independent calculation reproduced the response means exactly. The definitions in Table~\ref{tab:supp-native-definitions} are transcribed from the run configuration, not recomputed. The package holds the training mean, directions, random-direction seeds, edit sizes and the post hoc outcome bootstrap on available cases with per-draw values.
\item[$\pi_0$ rerun.] The records of the $\pi_0$ rerun on different hardware and of its admission rule are included (Sec.~\ref{sec:supp-pi0-internals}).
\end{description}
